Drug-resistant tuberculosis remains a major global health threat, yet the diagnostics in widest use are either too slow (culture-based drug susceptibility testing, 3--16 weeks) or too narrow (GeneXpert, rifampicin only), while rule-based whole-genome sequencing catalogues provide limited quantitative interpretability at the level of an individual isolate. Here we present FORUM-TB, an open and reproducible framework that converts public \textit{Mycobacterium tuberculosis} sequencing data into interpretable, multi-drug resistance predictions. We assembled a machine-learning-ready feature matrix from 9{,}798 publicly available isolates, encoding 2{,}693 nucleotide positions across nine genes associated with resistance to the four first-line drugs (rifampicin, isoniazid, ethambutol, and pyrazinamide), and benchmarked seven classifiers per drug alongside a multi-label formulation that predicts all four drugs jointly. Random Forest achieved the highest cross-validated AUC-ROC for every drug (rifampicin 0.969, isoniazid 0.946, ethambutol 0.900, pyrazinamide 0.883), consistently outperforming gradient-boosted and linear alternatives. Applying SHAP (SHapley Additive exPlanations) to the resulting models recovered the canonical resistance-determining mutations \textit{rpoB} S450L, \textit{katG} S315T, and \textit{embB} M306I/V as top-ranked features without these associations being supplied as prior knowledge, and a cross-drug attribution analysis quantified how the importance of shared features shifts under co-resistance. The multi-label extension reached a macro AUC-ROC of 0.919 but a markedly lower four-way exact-match accuracy of 0.591, indicating that joint multi-drug profiling is substantially harder than per-drug discrimination. These results show that an interpretable, per-drug ensemble built entirely from public data and open-source tooling recovers established resistance biology while surfacing quantitative co-resistance signal that static catalogues do not expose. The dataset, code, trained models, and an interactive dashboard are released openly to support reproducible benchmarking of interpretable resistance prediction.
